\input{../tex-inputs/latex-leseansicht-vorspann}

\section[Gustav Schwarzkopf u.a. an Arthur Schnitzler, 22. 7. 1904]{L04311 Gustav Schwarzkopf u.a. an Arthur Schnitzler, 22. 7. 1904}
\nopagebreak\mylabel{L04311v}
\rehead{ }\normalsize\beginnumbering\briefempfaengerindex{Schnitzler, Arthur@\textsc{Schnitzler, Arthur}!zzzSchönherr, Malvine@\emph{von Malvine Schönherr}!1904-07-221@{22. 7. 1904}|(be}\briefempfaengerindex{Schnitzler, Arthur@\textsc{Schnitzler, Arthur}!zzzChiavacci, Vincenz@\emph{von Vincenz Chiavacci}!1904-07-221@{22. 7. 1904}|(be}\briefempfaengerindex{Schnitzler, Arthur@\textsc{Schnitzler, Arthur}!zzzHirschfeld, Robert@\emph{von Robert Hirschfeld}!1904-07-221@{22. 7. 1904}|(be}\briefempfaengerindex{Schnitzler, Arthur@\textsc{Schnitzler, Arthur}!zzzHirschfeld, Emilie@\emph{von Emilie Hirschfeld}!1904-07-221@{22. 7. 1904}|(be}\briefempfaengerindex{Schnitzler, Arthur@\textsc{Schnitzler, Arthur}!zzzSchwarzkopf, Gustav@\emph{von Gustav Schwarzkopf}!1904-07-221@{22. 7. 1904}|(be}
\toendnotes[C]{\smallbreak\pagebreak[2]}
\correspDesc{Versand  durch Gustav Schwarzkopf, Emilie Hirschfeld, Robert Hirschfeld, Vinzenz Chiavacci, Malvine Chiavacci am 22. 7. 1904 in Sekirn
\newline{}Übermittlung  am 23. 7. 1904 in Klagenfurt
\newline{}Zustellung  am 25. 7. 1904 in Wien
\newline{}Erhalt  durch Arthur Schnitzler am 25. 7. 1904 in Wien}\toendnotes[C]{\smallbreak}
\Standort{CUL, Schnitzler, B 96a.}
\physDesc{Bildpostkarte, 246 Zeichen
\newline{}Handschrift Gustav Schwarzkopf: 1) Bleistift, deutsche Kurrent\hspace{1em}2) Bleistift, deutsche Kurrent\hspace{1em}
\newline{}Handschrift Emilie Hirschfeld: Bleistift, deutsche Kurrent
\newline{}Handschrift Robert Hirschfeld: Bleistift, deutsche Kurrent
\newline{}Handschrift Vincenz Chiavacci: Bleistift, deutsche Kurrent
\newline{}Handschrift Malvine Schönherr: Bleistift, deutsche Kurrent
\newline{}Versand: 1) Stempel: »\nobreak{}\oindex{Klagenfurt@\textbf{Klagenfurt}|pwk}Klagenfurt, 23 VII 04\nobreak{}«.   2) Stempel: »\nobreak{}\oindex{Sekirn@\textbf{Sekirn}|pwk}Sekirn\nobreak{}«.  3) Stempel: »\nobreak{}\oindex{XVIII., Währing@\textbf{XVIII., Währing}, \emph{Verwaltungsgebiet}|pwk}18/1 Wien 110, 25. 7. 04, bestellt\nobreak{}«. }\toendnotes[C]{\smallbreak}\pstart{}{\pb}Herrn \textsc{D\textsuperscript{r} Arthur Schnitzler}\pend{}\pstart{}\textsc{Wien\oindex{Wien@\textbf{Wien}, \emph{Verwaltungsgebiet}|pw}}\pend{}\pstart{}XVIII Spöttelgaſſe 7\oindex{Wien@\textbf{Wien}!XVIII., Währing@\textbf{XVIII., Währing}!Edmund-Weiß-Gasse 7@\textbf{Edmund-Weiß-Gasse 7}, \emph{Wohngebäude}|pw}.\pend{}{\bigskip}
\pstart
           \noindent{}\centering{}{\pb}\textcolor{gray}{\textbf{Jungbauer am Wörtherſee\oindex{Restauration Jungbauer am Wörthersee@\textbf{Restauration Jungbauer am Wörthersee}, \emph{Restaurant}|pw}.}}\pend
           \vspace{1em}
\pstart
           \raggedleft{}{\pb}22. VII. 04\pend
           \vspace{0.5em}
\pstart
           Sekirn\oindex{Sekirn@\textbf{Sekirn}|pw} \pend
           
\pstart
           Herzlichſt grüßt\pend
           \pstart \spacefill\mbox{Gustav Schwarzkopf}\pend{}\selectlanguage{ngerman}\vspace{1em}
\pstart
           \noindent{}{[}hs. Hirschfeld:{]} Ebenfalls\pend
           \pstart \spacefill\mbox{Hirschfeld}\pend{}\selectlanguage{ngerman}\vspace{1em}
\pstart
           \noindent{}{[}hs. Chiavacci:{]} Nicht minder\pend
           \pstart \spacefill\mbox{Chiavacci.}\pend{}\selectlanguage{ngerman}\vspace{1em}
\pstart
           \noindent{}{[}hs. Hirschfeld:{]} Ganz beſonders\pend
           \pstart \spacefill\mbox{Emmy Hirschfeld.}\pend{}\selectlanguage{ngerman}\vspace{1em}
\pstart
           \noindent{}{[}hs. Schönherr:{]} Herzliche Empfehlungen Ihner  liebenswürdigen Frau\pwindex{Schnitzler, Olga 17.\,1.\,1882 Wien – 13.\,1.\,1970 Lugano@\textsc{Schnitzler, Olga} (17.\,1.\,1882 Wien – 13.\,1.\,1970 Lugano), \emph{Schauspielerin, Sängerin}|pwv}\pend
           
\pstart
           \textcolor{gray}{herzlich}{\\[\baselineskip]}\spacefill\mbox{Malvine Chiavacci}\pend
           \leftskip=0em{}\selectlanguage{ngerman}\endnumbering\briefempfaengerindex{Schnitzler, Arthur@\textsc{Schnitzler, Arthur}!zzzSchönherr, Malvine@\emph{von Malvine Schönherr}!1904-07-221@{22. 7. 1904}|)be}\briefempfaengerindex{Schnitzler, Arthur@\textsc{Schnitzler, Arthur}!zzzChiavacci, Vincenz@\emph{von Vincenz Chiavacci}!1904-07-221@{22. 7. 1904}|)be}\briefempfaengerindex{Schnitzler, Arthur@\textsc{Schnitzler, Arthur}!zzzHirschfeld, Robert@\emph{von Robert Hirschfeld}!1904-07-221@{22. 7. 1904}|)be}\briefempfaengerindex{Schnitzler, Arthur@\textsc{Schnitzler, Arthur}!zzzHirschfeld, Emilie@\emph{von Emilie Hirschfeld}!1904-07-221@{22. 7. 1904}|)be}\briefempfaengerindex{Schnitzler, Arthur@\textsc{Schnitzler, Arthur}!zzzSchwarzkopf, Gustav@\emph{von Gustav Schwarzkopf}!1904-07-221@{22. 7. 1904}|)be}\mylabel{L04311h}
\begin{anhang}
\end{anhang}\newcommand{\dateiname}{L04311}\newcommand{\titel}{Gustav Schwarzkopf u.a. an Arthur Schnitzler, 22. 7. 1904}\newcommand{\editorInnen}{Herausgegeben von Jahnke, SelmaMüller, Martin Anton}\input{../tex-inputs/latex-leseansicht-abspann}
